% COLLECTION_METADATA: {"schema_version": 1, "kind": "reused_writeup", "independently_reviewed": false, "primary_source": "attacks/erdos/gpt_pro_5.2/51_v2.tex", "source_paths": ["attacks/erdos/gpt_5.2/51.tex", "attacks/erdos/gpt_pro_5.2/51.tex", "attacks/erdos/gpt_pro_5.2/51_v2.tex"], "source_model": "gpt_pro_5.2", "source_completion": 40.0, "source_claim": "unresolved"}
\section*{Erd\H{o}s Problem 51 --- collection record}

\subsection*{PROVENANCE AND STATUS}
Official problem: https://www.erdosproblems.com/51

Saved collaborative-database status: open.
Snapshot checked: 2026-09-23T17:22:24Z.
Latest status and formalization links: https://teorth.github.io/erdosproblems/

This is an attributed collection of existing writeups. The model-folder names below identify their original attribution; placement in GPT\_6\_Astra\_Ultra does not attribute their mathematical work to that model. All locally available source versions selected for this problem are reproduced below, without editing their text.

Proof claims, computations, literature statements, and completion estimates in the reused text have not been independently verified in this collection pass. A source claiming a solution does not change the saved upstream status.

Primary source (latest numbered version in the preferred model folder): \texttt{attacks/erdos/gpt\_pro\_5.2/51\_v2.tex}.

\subsection*{REUSED SOURCE 1}
Original attribution: \texttt{gpt\_5.2}.
Original file: \texttt{attacks/erdos/gpt\_5.2/51.tex}.

% BEGIN REUSED SOURCE: attacks/erdos/gpt_5.2/51.tex

OUTPUT
1) FORMAL RESTATEMENT
Definitions/Conventions.
- $\phi(n)$ is Euler's totient function.
- For $a\in\mathbb N$, say $a$ is a totient if $\exists n\in\mathbb N$ with $\phi(n)=a$.
- For a totient $a$, define
\[n_a:=\min\{n\in\mathbb N:\phi(n)=a\}.\]

Question.
Does there exist an infinite set $A\subset\mathbb N$ of totients such that
\[\lim_{a\to\infty,\ a\in A}\frac{n_a}{a}=+\infty?\]

2) QUICK LITERATURE/CONTEXT CHECK
Web browsing is not available in this session.
From the problem text only:
- Carmichael asked whether $\phi(n)=t$ can have exactly one solution.
- Erd\H{o}s proved that if such a $t$ exists then infinitely many such $t$ exist.
I do not use these claims as lemmas below.

3) ATTACK PLAN
Proof track ideas (existence of such an $A$):
1. Try to force any solution $n$ of $\phi(n)=a$ to have many distinct small prime factors, making $n/a=n/\phi(n)=\prod_{p\mid n}p/(p-1)$ very large.
2. Pick $a$ with carefully engineered divisibility constraints so that any preimage must include certain primes (e.g. by forcing many primes $p$ with $p-1\mid a$),
   and then argue that a smaller preimage cannot exist.

Disproof track ideas (nonexistence):
1. Try to show a universal upper bound $n_a \le C a (\log\log a)^{O(1)}$ or even $n_a\le Ca$ for all totients $a$.
2. Look for mechanisms producing small preimages, e.g. primes $p$ with $p-1=a$ or other constructions.

Chosen path in this attempt: prove two unconditional lemmas about the shape of totients/minimal preimages and run a sanity-check computation of $n_a/a$ for $a$ up to $3\times 10^5$.

4) WORK
Lemma 1 (Parity obstruction: $\phi(n)$ is even for $n>2$).
If $n>2$ then $\phi(n)$ is even. In particular, every totient $a>2$ is even.

Proof.
If $n$ is divisible by an odd prime $p$, then $\phi(n)$ is divisible by $\phi(p)=p-1$, which is even, so $\phi(n)$ is even.
If $n$ has no odd prime factor, then $n=2^k$ with $k\ge 2$ (since $n>2$), and $\phi(2^k)=2^{k-1}$ is even.
Thus in all cases $n>2\Rightarrow \phi(n)$ even. \qed

Lemma 2 (Prime totients have tiny minimal preimage).
Let $p$ be an odd prime and set $a:=p-1$. Then $a$ is a totient (since $\phi(p)=p-1$), and the minimal preimage equals $n_a=p$.
Consequently,
\[\frac{n_a}{a}=\frac{p}{p-1}\to 1\quad\text{as }p\to\infty.\]

Proof.
We have $\phi(p)=p-1=a$, so $n_a\le p$.
On the other hand, for any $n>1$, $\phi(n)\le n-1$ (since among $1,\dots,n$ only $n$ itself is not counted, and at least one number is not coprime unless $n$ is prime; the inequality $\phi(n)\le n-1$ is always true).
Thus if $\phi(n)=a=p-1$, then $n-1\ge \phi(n)=p-1$, so $n\ge p$.
Therefore $n_a=p$. The ratio statement is immediate. \qed

Lemma 3 (An infinite family with uniformly bounded $n_a/a$).
Let $m\ge 1$ and set $a:=2^m$. Then $a$ is a totient because $\phi(2^{m+1})=2^m=a$, and hence
\[\frac{n_a}{a}\le \frac{2^{m+1}}{2^m}=2.\]

Proof.
Compute $\phi(2^{m+1})=2^{m+1}-2^m=2^m=a$. By definition of $n_a$ as the smallest preimage, $n_a\le 2^{m+1}$.
Divide by $a=2^m$ to get $n_a/a\le 2$. \qed

FAST REALITY CHECK (local computation; $a\le 300{,}000$).
I computed $\phi(n)$ for $1\le n\le 300{,}000$ and for each totient value $a\le 300{,}000$ recorded the smallest $n_a\le 300{,}000$ with $\phi(n_a)=a$.
Among these $a$, the largest ratios $n_a/a$ observed were about $2.0$:
top examples (within this search bound):
\[
\begin{array}{c|c|c}
a & n_a & n_a/a\\\hline
5888 & 11985 & 2.035\\
10496 & 21165 & 2.016\\
17408 & 34935 & 2.007\\
32768 & 65535 & 2.000
\end{array}
\]
Among $a\ge 50{,}000$ (still with $n_a\le 300{,}000$), the largest ratios found were $\approx 1.996$.
This is consistent with (but does not prove) the possibility that $n_a/a$ might remain bounded, contrary to the question's desired conclusion.

5) VERIFICATION
- Lemma 1 is standard and the proof splits into ``odd prime factor'' vs ``power of two'' cases.
- Lemma 2 uses only the elementary inequality $\phi(n)\le n-1$ and the explicit solution $n=p$.
- Lemma 3 is a direct computation.
- Computation is limited to preimages $n\le 300{,}000$, so it cannot rule out larger minimal preimages for larger $a$.

6) FINAL
**UNRESOLVED**

(i) Strongest fully proved partial result:
I proved basic obstructions and easy families:
every totient $a>2$ is even (Lemma 1),
totients of the form $a=p-1$ (odd prime $p$) have $n_a/a\to 1$ (Lemma 2),
and the infinite family $a=2^m$ has $n_a/a\le 2$ (Lemma 3).

(ii) First gap:
Construct an infinite set $A$ of totients for which \emph{every} solution $n$ to $\phi(n)=a$ is forced to have very large $n/a$,
or prove a general upper bound on $n_a/a$ that rules out divergence.

(iii) Top 3 next moves (concrete):
1. Try to force many primes into any preimage: choose $a$ so that for many primes $p$ one has $p-1\mid a$, and show that any $n$ with $\phi(n)=a$ must be divisible by many of those $p$.
2. Study the structure of minimal preimages: prove lemmas of the form ``if $\phi(n)=a$ then there exists a solution $n'\le C(a)$ with specified constraints'', aiming to bound $n_a$.
3. Extend computation to much larger bounds (e.g. $n\le 10^8$) and track the growth of $\max_{a\le X} n_a/a$; if it stays bounded, attempt to conjecture a sharp constant.

(iv) What a minimal counterexample would likely look like (structure/parameters):
An infinite sequence of totients $a_k$ such that every $n$ with $\phi(n)=a_k$ must include an increasing set of small primes in its factorization,
forcing $n/\phi(n)=\prod_{p\mid n}p/(p-1)\to\infty$, while simultaneously preventing the existence of alternative preimages with fewer primes (and hence smaller $n$).


% END REUSED SOURCE: attacks/erdos/gpt_5.2/51.tex

\subsection*{REUSED SOURCE 2}
Original attribution: \texttt{gpt\_pro\_5.2}.
Original file: \texttt{attacks/erdos/gpt\_pro\_5.2/51.tex}.

% BEGIN REUSED SOURCE: attacks/erdos/gpt_pro_5.2/51.tex

\subsection*{Erd\H{o}s Problem \#51}
\noindent\emph{Problem statement (from the problem file).}
Is there an infinite set $A\subset\mathbb{N}$ such that for every $a\in A$ there is an integer $n$ with $\varphi(n)=a$, and yet if $n_a$ is the smallest such integer then $n_a/a\to\infty$ as $a\to\infty$?

\paragraph{FORMAL RESTATEMENT.}
Does there exist an infinite set $\mathcal{A}\subset\mathbb{N}$ such that
\begin{itemize}
\item for every $a\in\mathcal{A}$, the equation $\varphi(n)=a$ has at least one solution $n\in\mathbb{N}$;
\item writing $n_a:=\min\{n\in\mathbb{N}:\varphi(n)=a\}$, one has $\lim_{\substack{a\to\infty\\ a\in\mathcal{A}}} \frac{n_a}{a}=\infty$.
\end{itemize}

\paragraph{QUICK LITERATURE/CONTEXT CHECK.}
The problem file mentions Carmichael's question about uniqueness of solutions to $\varphi(n)=t$ and that Erd\H{o}s proved: if there exists such a $t$ with a unique preimage then there are infinitely many.
I do not use any further results.

\paragraph{ATTACK PLAN.}
\emph{Construction track:} Search for a family of totients $a$ that force the smallest preimage to have many small prime factors, making $n_a/a$ large.
\emph{Disproof track:} Try to show that for every totient value $a$ there is always some preimage $n$ with $n\ll a(\log\log a)$ or similar, which would preclude divergence.
\emph{Reality check:} Compute $n_a$ for all $a$ realized by $\varphi(n)$ up to a cutoff and inspect ratios.

\paragraph{WORK.}
\subparagraph{FAST REALITY CHECK (numerics).}
Scanning all $n\le 300000$ and recording, for each totient value $a=\varphi(n)\le 20000$, the minimal preimage $n_a$, we found:
\begin{itemize}
\item among the $4486$ values $a\le 20000$ that occur as $\varphi(n)$ with some $n\le 300000$, the maximum observed ratio $n_a/a$ was about $2.035$ (attained at $a=5888$, $n_a=11985$).
\item for values of the form $a=p-1$ with $p$ prime, one has $n_a=p$ and $n_a/a=1+1/a$, quickly tending to $1$.
\end{itemize}
This finite search does not show large ratios, but it is far from the asymptotic regime.

\subparagraph{Lemma 1 (basic ratio identity).}
For every $n\ge 1$,
\[
\frac{n}{\varphi(n)}=\prod_{p\mid n}\frac{1}{1-1/p}.
\]

\emph{Proof.}
This is immediate from Lemma~1 of Problem~\#50: $\varphi(n)/n=\prod_{p\mid n}(1-1/p)$, hence invert. \qed

\subparagraph{Lemma 2 (prime-shifted totients have minimal preimage equal to the prime).}
If $p$ is prime and $a:=p-1$, then the smallest $n$ with $\varphi(n)=a$ is $n_a=p$.

\emph{Proof.}
We have $\varphi(p)=p-1=a$, so $n_a\le p$.
On the other hand, for every integer $n>1$, one has $\varphi(n)\le n-1$ with equality only when $n$ is prime.
Thus if $\varphi(n)=p-1$, then $n\ge p$.
Combining gives $n_a=p$. \qed

\subparagraph{Lemma 3 (a general lower bound).}
For any $a\ge 1$ that occurs as $a=\varphi(n)$ for some $n$, one has $n_a\ge a+1$.

\emph{Proof.}
If $a=1$ then $n_a=1$ and the claim holds.
For $a\ge 2$, any $n$ with $\varphi(n)=a$ satisfies $n>1$ and hence $\varphi(n)\le n-1$.
Thus $a\le n-1$, i.e. $n\ge a+1$.
Taking the minimum over solutions gives $n_a\ge a+1$. \qed

\paragraph{VERIFICATION.}
\begin{itemize}
\item Lemma~2 uses only the elementary inequality $\varphi(n)\le n-1$ for $n>1$.
\item Lemma~1 is the standard prime-factor identity; it shows how large ratios require many small prime factors.
\item The numerical scan is explicitly bounded and cannot resolve asymptotics.
\end{itemize}

\paragraph{FINAL.} \textbf{UNRESOLVED.}
\begin{enumerate}
\item[(i)] \emph{Strongest proved partial result.} The ratio $n/\varphi(n)$ is the multiplicative product over primes dividing $n$ (Lemma~1), and for $a=p-1$ one has $n_a=p$ so $n_a/a\to 1$ along this infinite family (Lemma~2). Any totient value $a$ satisfies $n_a\ge a+1$ (Lemma~3).
\item[(ii)] \emph{First gap.} Construct an explicit infinite set $\mathcal{A}$ of totient values for which the minimal preimage $n_a$ necessarily has $n_a/a\to\infty$, or prove that such a set cannot exist.
\item[(iii)] \emph{Top 3 next moves.}
(1) Seek a family of totient values $a$ whose preimages must include many small primes (forcing large $n/\varphi(n)$), and then prove that no smaller $n$ can realize the same $a$.
(2) Experiment computationally with larger cutoffs to see whether $\max_{a\le X} n_a/a$ grows, and if so for which structural $a$ (highly composite, smooth, etc.).
(3) Prove structural lemmas about the minimal preimage: e.g. constraints on its prime factorization or on the largest prime factor in terms of $a$.
\item[(iv)] \emph{Minimal counterexample structure.} If the answer is ``no'', one would expect a universal function $g(a)=o(\infty)$ such that for every totient value $a$ there exists some $n\le g(a)\,a$ with $\varphi(n)=a$; a counterexample to this would be a sequence of totients $a_k$ whose all preimages satisfy $n/a_k\to\infty$.
\end{enumerate}


% END REUSED SOURCE: attacks/erdos/gpt_pro_5.2/51.tex

\subsection*{REUSED SOURCE 3}
Original attribution: \texttt{gpt\_pro\_5.2}.
Original file: \texttt{attacks/erdos/gpt\_pro\_5.2/51\_v2.tex}.

% BEGIN REUSED SOURCE: attacks/erdos/gpt_pro_5.2/51_v2.tex
\section{Erd\H{o}s Problem \#51 --- Round 2}

\subsection*{1) ROUND-2 OBJECTIVE}
\textbf{Path (C): obstruction/bounding approach.}
Round~1 identified that large values of $n_a/a$ require $n_a$ to have many small prime divisors (Lemma~1 there), but left open whether this can be forced for the \emph{minimal} preimage.
The most promising next step is therefore to (i) obtain \emph{unconditional} global upper bounds on $n_a/a$ (ruling out any ``fast'' divergence), and (ii) push computation far enough to see whether record values of $n_a/a$ grow at all.

\subsection*{2) ROUND-1 FOUNDATION USED}
We rely on the following Round~1 items (without reproving them):
\begin{itemize}
\item (R1--Lemma~1) $\displaystyle \frac{n}{\varphi(n)}=\prod_{p\mid n}\frac{p}{p-1}$.
\item (R1--Lemma~2) If $a=p-1$ with $p$ prime then $n_a=p$ and $n_a/a\to 1$ along this infinite family.
\item (R1--Lemma~3) For any totient value $a$, $n_a\ge a+1$.
\item (R1--Computation) For $a\le 20000$ the maximum observed ratio was $\approx 2.035$.
\end{itemize}

\subsection*{3) NEW INSIGHT / TOOL (ROUND-2)}
\begin{itemize}
\item \textbf{New structural lemma:} $\varphi(n)\ge \sqrt n$ for all $n\neq 2,6$. This implies \emph{quadratic compression} for inverse totients: if $\varphi(n)=a$ and $a\ge 3$ then $n\le a^2$, hence $n_a\le a^2$.
\item \textbf{New explicit analytic bound:} using Rosser--Schoenfeld's inequality
\[\frac{n}{\varphi(n)}<e^{\gamma}\log\log n+\frac{2.50637}{\log\log n}\qquad(n\ge 3),\]
we obtain an explicit upper bound on $n_a/a$ in terms of $a$ alone.
\item \textbf{Much larger computation:} a totient sieve up to $n\le 10^8$ to locate record values of $n_a/a$.
\end{itemize}

\subsection*{4) ATTACK PLAN (ROUND-2)}
\begin{enumerate}
\item Close a Round~1 gap by proving a nontrivial \emph{unconditional} upper bound for $n_a/a$.
\item Extend the computational search from $n\le 3\cdot 10^5$ (Round~1) to $n\le 10^8$, recording the maximum of $n_a/a$ among all totients realized with minimal preimage $\le 10^8$.
\item Try to identify rigid structural patterns among record-holders as a guide for (future) proofs or counterexamples.
\end{enumerate}

\subsection*{5) WORK (ROUND-2)}

\subsubsection*{5.1. A basic but powerful inequality: $\varphi(n)\ge \sqrt n$ except $n=2,6$}
\begin{lemma}[\,$\varphi(n)\ge \sqrt n$ for $n\neq 2,6$\,]
For every integer $n\ge 1$ with $n\neq 2,6$, we have $\varphi(n)\ge \sqrt n$.
Equivalently, $\varphi(n)^2\ge n$ for all $n\neq 2,6$.
\end{lemma}

\begin{proof}
Define
\[g(n):=\frac{\varphi(n)^2}{n}.
\]
Since $\varphi$ is multiplicative and $n\mapsto 1/n$ is multiplicative, $g$ is multiplicative:
if $(m,n)=1$ then
\[
 g(mn)=\frac{\varphi(mn)^2}{mn}=\frac{\varphi(m)^2}{m}\cdot\frac{\varphi(n)^2}{n}=g(m)g(n).
\]
Hence it suffices to analyze prime powers.
For $p$ prime and $k\ge 1$,
\[
 g(p^k)=\frac{\varphi(p^k)^2}{p^k}=\frac{\bigl(p^{k-1}(p-1)\bigr)^2}{p^k}=p^{k-2}(p-1)^2.
\]
\emph{Odd primes.}
If $p\ge 3$ then for all $k\ge 1$,
\[
 g(p^k)\ge g(p)=\frac{(p-1)^2}{p} > 1,
\]
since $(p-1)^2-p=p^2-3p+1\ge 1$ for $p\ge 3$.
So $g(p^k)>1$ for every odd prime power.

\emph{The prime $2$.}
We have $g(2)=\frac{1}{2}<1$, while for $k\ge 2$,
\[
 g(2^k)=2^{k-2}\ge 1.
\]
Now write $n=2^k m$ with $m$ odd.
If $k\ge 2$ then $g(n)=g(2^k)g(m)\ge 1\cdot 1=1$.
If $k=0$ then $n$ is odd and $g(n)=g(m)\ge 1$ (with equality only at $m=1$).
If $k=1$ then $g(n)=g(2)g(m)=\tfrac12 g(m)$.
This is $<1$ only if $g(m)<2$.
But for odd $m>1$, the smallest possible value of $g(m)$ occurs at $m=3$, giving $g(3)=4/3<2$; for $m\ge 5$, one checks $g(m)\ge g(5)=16/5>2$ because every prime-power factor contributes $g(p^k)>1$ and the least such factor for odd primes is $g(5)=16/5$.
Hence $g(2m)<1$ occurs exactly for $m=1$ (i.e. $n=2$) and $m=3$ (i.e. $n=6$).

Therefore $g(n)\ge 1$ for all $n\neq 2,6$, i.e. $\varphi(n)^2\ge n$ and $\varphi(n)\ge \sqrt n$ for $n\neq 2,6$.
\end{proof}

\begin{corollary}[Quadratic bound for inverse totients]
If $a\ge 3$ is a totient value and $n$ is any solution of $\varphi(n)=a$, then $n\le a^2$.
In particular, $n_a\le a^2$ for all totient values $a\ge 3$.
\end{corollary}

\begin{proof}
Let $\varphi(n)=a$ with $a\ge 3$.
Then $n\neq 2,6$ (since $\varphi(2)=1$ and $\varphi(6)=2$), so by the lemma $n\le \varphi(n)^2=a^2$.
Taking the minimum over solutions gives $n_a\le a^2$.
\end{proof}

\subsubsection*{5.2. An explicit analytic upper bound on $n_a/a$}
We use the following explicit inequality of Rosser--Schoenfeld.

\begin{theorem}[Rosser--Schoenfeld (1962), explicit upper bound]
For every integer $n\ge 3$,
\[
\frac{n}{\varphi(n)}<e^{\gamma}\log\log n+\frac{2.50637}{\log\log n}.
\]
\end{theorem}

\begin{theorem}[Explicit upper bound for $n_a/a$]
Let $a\ge 3$ be a totient value and $n_a:=\min\{n:\varphi(n)=a\}$.
Then
\[
\frac{n_a}{a}<e^{\gamma}\log\log(a^2)+\frac{2.50637}{\log\log(a^2)}
\;=\;
 e^{\gamma}\bigl(\log\log a+\log 2\bigr)+\frac{2.50637}{\log\log a+\log 2}.
\]
In particular, $\frac{n_a}{a}=O(\log\log a)$.
\end{theorem}

\begin{proof}
Since $a\ge 3$, we have $n_a\ge a+1\ge 4$, hence $n_a\ge 3$.
Applying Rosser--Schoenfeld to $n=n_a$ gives
\[
\frac{n_a}{a}=\frac{n_a}{\varphi(n_a)}<e^{\gamma}\log\log n_a+\frac{2.50637}{\log\log n_a}.
\]
By the previous corollary, $n_a\le a^2$.
As $\log\log$ is increasing on $[3,\infty)$, we have $\log\log n_a\le \log\log(a^2)$.
Substituting yields the displayed bound.
The $O(\log\log a)$ conclusion follows immediately.
\end{proof}

\subsubsection*{5.3. Extended computation up to $n\le 10^8$}
\paragraph{Method.}
We computed $\varphi(n)$ for all $1\le n\le 10^8$ by a standard totient sieve.
For each totient value $a\le 10^8$ encountered, we recorded the least $n$ with $\varphi(n)=a$.
Finally we scanned over all such $a$ to find the maximum of $n_a/a$.

\paragraph{Result (record up to $10^8$).}
Among all totients $a$ with minimal preimage $n_a\le 10^8$, the maximum observed ratio was
\[
\max \frac{n_a}{a} \approx 2.041378860109,
\]
attained at
\[
 a=2037248,\qquad n_a=4158795=3\cdot 5\cdot 17\cdot 47\cdot 347.
\]
For this record value,
\[
\frac{n_a}{a}=\frac{4158795}{2037248}=\prod_{p\mid n_a}\frac{p}{p-1}
 =\frac32\cdot\frac54\cdot\frac{17}{16}\cdot\frac{47}{46}\cdot\frac{347}{346}.
\]

\paragraph{Top record-holders (up to $10^8$).}
The next-largest ratios observed are listed below.
Empirically, every entry has the form $n_a=3\cdot 5\cdot 17\cdot 47\cdot q$ with $q$ prime, and hence
$a=\varphi(n_a)=2\cdot 4\cdot 16\cdot 46\cdot (q-1)=5888\,(q-1)$.
\begin{center}
\begin{tabular}{r|r|r|l}
rank & $a$ & $n_a$ & $n_a/a$ \\\hline
1 & 2037248 & 4158795 & 2.0413788601\\
2 & 6983168 & 14226195 & 2.0372121937\\
3 & 8043008 & 16383495 & 2.0369860381\\
4 & 8749568 & 17821695 & 2.0368657058\\
5 & 11929088 & 24293595 & 2.0365006109\\
6 & 12423680 & 25300335 & 2.0364606139\\
7 & 14578688 & 29686845 & 2.0363180144\\
8 & 23587328 & 48023895 & 2.0360040357\\
9 & 25212416 & 51331755 & 2.0359712849\\
10 & 26236928 & 53417145 & 2.0359527228\\
\end{tabular}
\end{center}

\subsection*{6) ADVERSARIAL VERIFICATION}
\begin{itemize}
\item \textbf{Lemma $\varphi(n)\ge \sqrt n$ (except $2,6$).}
We reduced to prime powers via the multiplicativity of $g(n)=\varphi(n)^2/n$.
All prime-power factors satisfy $g(p^k)>1$ except $g(2)=1/2$.
The only way for a product of such factors to fall below $1$ is to include $g(2)$ and multiply it by an odd factor $g(m)<2$; checking odd $m$ shows only $m=1,3$ work, yielding $n=2,6$.
Boundary checks: $\varphi(2)=1<\sqrt2$ and $\varphi(6)=2<\sqrt6$ do fail, as required.
\item \textbf{Rosser--Schoenfeld bound usage.}
We only apply it for $n=n_a\ge 4$, so $\log\log n_a>0$.
The step $\log\log n_a\le \log\log(a^2)$ uses the proven inequality $n_a\le a^2$.
\item \textbf{Computation.}
The record value $n_a=3\cdot 5\cdot 17\cdot 47\cdot 347$ was independently verified by direct factorization and by checking $\varphi(n_a)=a$ via the product formula.
The computation is still finite: it does not exclude larger ratios occurring at $n_a>10^8$.
\end{itemize}

\subsection*{7) FINAL}
\textbf{UNRESOLVED (BUT STRICTLY ADVANCED).}
We proved nontrivial global structure and bounds:
\begin{itemize}
\item $\varphi(n)\ge \sqrt n$ for $n\neq 2,6$, implying $n_a\le a^2$ for every totient $a\ge 3$.
\item Using Rosser--Schoenfeld, we derived an explicit unconditional upper bound
$\displaystyle \frac{n_a}{a}<e^{\gamma}\log\log(a^2)+\frac{2.50637}{\log\log(a^2)}$, hence $n_a/a=O(\log\log a)$.
\item Computation up to $n_a\le 10^8$ finds a \emph{stable} maximum ratio $\approx 2.0413788601$ (no growth beyond the Round~1 range), with record-holders of a very rigid empirical shape.
\end{itemize}
The main Erd\H{o}s question (existence of an infinite set with $n_a/a\to\infty$) remains open.

\subsection*{8) COMPLETION ESTIMATE (MANDATORY)}
\textbf{COMPLETION: 40\%}

\subsection*{9) REFERENCES}
\begin{itemize}
\item J. B. Rosser and L. Schoenfeld, ``Approximate formulas for some functions of prime numbers'', \emph{Illinois J. Math.} \textbf{6} (1962), Theorem~15 (explicit bounds for $n/\varphi(n)$).
\item E. W. Weisstein, ``Totient Function'', MathWorld (states $\varphi(n)\ge \sqrt n$ for $n\neq 2,6$ and cites Kendall--Osborn (1965) and Mitrinovi\'c--S\'andor (1995)).
\end{itemize}

% END REUSED SOURCE: attacks/erdos/gpt_pro_5.2/51_v2.tex

\subsection*{COMPLETION ESTIMATE}
Inherited primary-source estimate: $40\%$. This is the original source's unverified estimate, not a new assessment or confirmation that the problem is solved.
